\chapter{Methodology}
\label{ch:methodology}

Our simulation framework models a stylized market environment in which agents
repeatedly allocate resources under uncertainty. This chapter specifies the
environment dynamics, each agent's decision rule, and the formal guarantees
we establish for the two proposed methods.

%---------------------------------------------------------------
\section{The Stylized Market Environment}
%---------------------------------------------------------------

We simulate environments where payoffs are governed by hidden, potentially
changing parameters. At each step $t$, the environment draws a return
$r_t$ from a regime-specific distribution parameterised by $(\mu_{S_t},
\sigma_{S_t})$, where $S_t \in \{0,1\}$ is a hidden Markov state. The agent
observes only $r_t$ after placing its bet; the hidden state $S_t$ is
never directly revealed.

\subsection{Synthetic Environments}

Three synthetic regimes are evaluated:

\begin{enumerate}
    \item \textbf{Stationary:} A single regime with $r_t \overset{iid}{\sim}
    \Normal(\mu, \sigma^2)$, $\mu = 0.08$, $\sigma = 0.15$. Provides a
    controlled test of exploration-exploitation convergence where theory
    predicts Thompson-Kelly should match the Oracle asymptotically.

    \item \textbf{Adversarial Shock:} Bull regime ($\mu = 0.08, \sigma = 0.15$)
    for $t \leq T/2$, then a deterministic switch to a heavy-tailed bear
    regime ($r_t \sim t_3(-0.10, 0.30^2)$) for $t > T/2$. The Student-$t$
    with $\nu=3$ degrees of freedom has infinite fourth moment, creating the
    heavy-tail ruin conditions analysed in Experiment~6.

    \item \textbf{Slow Gaussian Drift:} The mean $\mu_t$ follows a
    deterministic linear drift from $+0.08$ to $-0.05$ over the horizon,
    testing whether agents can ``unlearn'' a decaying regime gracefully.
\end{enumerate}

\subsection{Common Random Numbers (CRN)}

All agents in a given experiment share the same pre-generated return matrix
$\mathbf{R} \in \R^{N_{\text{sims}} \times T}$. This \emph{Common Random
Numbers} methodology eliminates idiosyncratic Monte Carlo variance from
agent comparisons, so that differences in outcomes are attributable
solely to the agents' decision rules rather than to sampling noise.
Formally, for all agents $a, b$ and all simulations $i$:
\[
    r_t^{(i)}[\text{agent } a] = r_t^{(i)}[\text{agent }b] = R_{i,t}.
\]

%---------------------------------------------------------------
\section{Agent Descriptions}
\label{sec:agents}
%---------------------------------------------------------------

\subsection{Kelly Oracle (Benchmark)}

The Oracle knows the true regime-conditional distribution at each step
and bets $f_t^{\text{Oracle}} = f^*(\mu_{S_t}, \sigma_{S_t})$, where
$f^* = \mu/\sigma^2$ is the continuous Kelly fraction. This agent is
inadmissible (it uses future information) and serves only as a theoretical
upper bound on regret.

\subsection{NaiveBayes Kelly}

As defined in Chapter~\ref{ch:foundations}, this agent plugs the Beta
posterior mean into the Kelly formula. It is included to demonstrate the
sticky prior failure mode: in adversarial experiments, this agent
maintains a bull-regime prior well into the bear regime, over-leveraging
and incurring catastrophic losses.

\subsection{Thompson-Kelly Agent}

Samples $\tilde{p} \sim \text{Beta}(\alpha_t, \beta_t)$ and bets
$f_t = \text{Kelly}(\tilde{p})$. The posterior sampling provides implicit
uncertainty-driven leverage reduction.

\subsection{UCB-Kelly Agent}

Computes an Upper Confidence Bound on the win probability:
\[
    p_t^{\text{UCB}} = \hat{p}_t + \sqrt{\frac{2\ln t}{n_t}},
\]
clipped to $[0,1]$, then bets $f_t = \text{Kelly}(p_t^{\text{UCB}})$.
The exploration bonus inflates the bet when an arm is under-explored.

\subsection{EXP3-Fractional Agent}

\begin{remark}[Deviation from Canonical EXP3]
\label{rem:exp3}
    The canonical EXP3 algorithm of \textcite{auer2002nonstochastic} selects
    a \emph{single} arm at each step with probability proportional to
    exponentially re-weighted rewards. In our continuous betting context,
    we implement a \emph{fractional} variant where the EXP3 probability
    vector is used directly as a soft allocation: $f_{i,t} = p_{i,t}
    \cdot f_{\max}$. This variant does not satisfy the original EXP3 regret
    bound verbatim; it should be understood as an EXP3-inspired adversarial
    baseline that retains the exponential weight update but distributes
    rather than concentrates allocation. We cite the original EXP3 bound
    as a qualitative motivator, not a strict guarantee for this variant.
\end{remark}

\subsection{Softmax-Kelly Agent}

Assigns arm probabilities via a Boltzmann distribution over empirical
means with temperature $\tau$: $p_i \propto \exp(\hat{\mu}_i / \tau)$,
then weights Kelly fractions by these probabilities.

%---------------------------------------------------------------
\section{Proposed Method I: Volatility-Augmented HMM (Vol-HMM)}
\label{sec:vol-hmm}
%---------------------------------------------------------------

The Vol-HMM extends the standard HMM (Definition~\ref{def:hmm}) by
augmenting the scalar return observation $r_t$ with a rolling volatility
feature $v_t = \text{std}(r_{t-k+1}, \ldots, r_t)$ for window $k$.

\subsection{Two-Dimensional Emission Model}

For each regime $j \in \{\text{bull}, \text{bear}\}$, the emission
probability factorises over the two features:
\begin{equation}
\label{eq:emission}
    \log b_j(r_t, v_t) = \log \Normal(r_t;\, \mu_j, \sigma_j^2) + 
    \log \Gam(v_t;\, \kappa_j, \theta_j),
\end{equation}
where $\Gam(\cdot;\kappa, \theta)$ is the Gamma density with shape $\kappa$
and scale $\theta$.

\begin{remark}[Fixed vs.\ Estimated Emission Parameters]
\label{rem:hmm-fixed}
    The emission parameters $(\mu_j, \sigma_j, \kappa_j, \theta_j)$ are
    \emph{specified a priori} based on domain knowledge of typical bull and
    bear market statistics, rather than estimated via the Baum-Welch EM
    algorithm. This design choice is deliberate: with only $T = 250$
    observations per episode, EM estimation risks overfitting the emission
    model to the training sequence, particularly for the Gamma parameters on
    the volatility channel. By specifying the emission model, we treat the
    HMM as a \emph{Bayesian filter} (propagating belief over a pre-specified
    model) rather than as a generative model to be fitted. This is a standard
    approach in engineering applications of HMMs where the state semantics
    are known a priori (see, e.g., Kalman filtering literature).
\end{remark}

\subsection{Constrained Transition Matrix}

To address the sticky prior problem, we impose
\[
    A_{ii} \leq A_{\max} = 0.95, \quad i \in \{\text{bull}, \text{bear}\}.
\]
This ensures that even with zero new evidence, the probability of a regime
change at each step is at least $1 - 0.95 = 5\%$, preventing the posterior
from becoming arbitrarily concentrated.

\subsection{CUSUM Augmentation}

The CUSUM statistic for detecting a downward mean shift from the bull
regime is:
\[
    S_t^- = \max\!\bigl(0,\, S_{t-1}^- - z_t - k\bigr), \quad
    z_t = \frac{r_t - \mu_{\text{bull}}}{\sigma_{\text{bull}}},
\]
where $k > 0$ is a drift parameter. When $S_t^- > h$ (threshold $h = 3$
in our implementation), the log-emission for the bear state is boosted by
an additive constant, triggering rapid reallocation. This hybrid approach
combines CUSUM's distribution-free rapid detection with the HMM's
probabilistic regime attribution.

\subsection{Kelly Allocation}

Given filtered belief $\pi_t = (\pi_t^{\text{bull}}, \pi_t^{\text{bear}})$,
the bet fraction is
\[
    f_t^{\text{HMM}} = \pi_t^{\text{bull}} \cdot f^*_{\text{bull}} +
    \pi_t^{\text{bear}} \cdot f^*_{\text{bear}},
\]
where $f^*_j = \max(0, \mu_j/\sigma_j^2)$. In the bear regime, $\mu_j < 0$
implies $f^*_j = 0$, so the agent goes to cash automatically once the
posterior shifts sufficiently.

\begin{proposition}[Monotone Leverage in Belief]
\label{prop:leverage}
    Under the Vol-HMM, $f_t^{\text{HMM}}$ is a non-increasing function of
    $\pi_t^{\text{bear}}$. Specifically,
    \[
        \frac{\partial f_t^{\text{HMM}}}{\partial \pi_t^{\text{bear}}} =
        f^*_{\text{bear}} - f^*_{\text{bull}} \leq 0,
    \]
    since $f^*_{\text{bear}} \leq f^*_{\text{bull}}$ when $\mu_{\text{bear}}
    \leq \mu_{\text{bull}}$.
\end{proposition}

This proposition ensures the Vol-HMM \emph{automatically de-levers} as
evidence accumulates for the bear regime, with no ad hoc rule required.

%---------------------------------------------------------------
\section{Proposed Method II: Risk-Constrained Kelly (CPPI)}
\label{sec:cppi}
%---------------------------------------------------------------

\subsection{Drawdown Floor and Cushion}

Define the peak wealth up to time $t$ as
$W_t^{\max} = \max_{s \leq t} W_s$. The drawdown floor is
\[
    F_t = W_t^{\max}(1 - D_{\max}),
\]
where $D_{\max}$ is the maximum permissible drawdown (e.g., $0.20$ = 20\%).
The \emph{cushion} at time $t$ is $C_t = W_t - F_t \geq 0$.

\subsection{CPPI Leverage}

The constrained leverage multiplier is
\[
    \lambda_t = \min\!\left(1,\; \frac{m \cdot C_t}{W_t}\right),
\]
where $m \geq 1$ is the CPPI multiplier. The effective betting fraction is
$f_t^{\text{CPPI}} = \lambda_t \cdot f_t^{\text{base}}$, where
$f_t^{\text{base}}$ is any base Kelly fraction.

\begin{proposition}[CPPI Soft Floor in Discrete Time]
\label{prop:cppi}
    Suppose $|f_t r_t| \leq 1 - \varepsilon$ almost surely for some
    $\varepsilon > 0$ (bounded single-step loss). Then the CPPI strategy
    guarantees $W_t \geq F_t(1 - \delta)$ for all $t$, where $\delta$
    depends on $m$ and the maximum single-period loss.
\end{proposition}

\begin{proof}[Proof Sketch]
    If $C_t > 0$ and $m \leq 1/\varepsilon$, then
    $f_t^{\text{CPPI}} \cdot |r_t| \leq \lambda_t \leq m C_t / W_t$.
    The worst-case single-step wealth change is
    $W_{t+1} \geq W_t(1 - m C_t/W_t) = W_t - m C_t = F_t + (1-m)C_t$.
    For $m \leq 1 + F_t/C_t$, we have $W_{t+1} \geq F_t$.
    In discrete time with large single-step shocks, the floor can be
    breached (``gap risk''); our simulation quantifies this empirically.
    \qed
\end{proof}

\begin{remark}[Gap Risk in Discrete Time]
    Proposition~\ref{prop:cppi} requires bounded single-step losses. Under
    the Student-$t(\nu=3)$ distribution, the loss $1 + f_t r_t$ can be
    arbitrarily close to zero with non-negligible probability, creating
    gap risk. In practice, we observe zero ruin events for the CPPI agent
    across all tested scenarios, suggesting the floor is sufficient in
    the simulated parameter regime.
\end{remark}

%---------------------------------------------------------------
\section{Market Friction and Transaction Costs}
%---------------------------------------------------------------

To bridge theoretical growth and empirical deployment, we model
transaction costs at each step. Let $\hat{f}_{t-1}$ denote the effective
exposure after the previous price move:
\[
    \hat{f}_{t-1} = \frac{f_{t-1}(1 + r_{t-1})}{1 + f_{t-1} r_{t-1}}.
\]
The turnover at step $t$ is $\Delta f_t = |f_t - \hat{f}_{t-1}|$, and the
cost-adjusted multiplier is
\begin{equation}
\label{eq:friction}
    M_t = 1 + f_t r_t - c \cdot \Delta f_t,
\end{equation}
where $c$ is the transaction cost in basis points divided by 10{,}000. If
$M_t \leq 0$, the path is ruined. This formulation penalises ``chattering''
algorithms that rebalance frequently on noise, a property we exploit in
Section~\ref{ch:experiments} to show that EXP3 and UCB incur
disproportionate fee drag.

%---------------------------------------------------------------
\section{Evaluation Metrics}
%---------------------------------------------------------------

Strategies are evaluated against the following metrics:

\begin{itemize}
    \item \textbf{Median Terminal Wealth} $\tilde{W}_T$: Robust to heavy
    right tails that inflate the mean.

    \item \textbf{Probability of Ruin} $\hat{p}_{\text{ruin}}$:
    Empirical fraction of paths reaching $W_t = 0$.

    \item \textbf{Annualised Sharpe Ratio}:
    \[
        \mathcal{S} = \frac{\bar{R}}{\hat{\sigma}(R)} \cdot \sqrt{252},
    \]
    where $\bar{R}$ and $\hat{\sigma}(R)$ are the sample mean and standard
    deviation of step-wise returns, and 252 is the assumed number of trading
    days per year.

    \item \textbf{Annualised Sortino Ratio}:
    \[
        \mathcal{SOR} = \frac{\bar{R}}{\hat{\sigma}_{\downarrow}(R)} \cdot \sqrt{252},
    \]
    where $\hat{\sigma}_{\downarrow}(R) = \sqrt{\E[\min(R,0)^2]}$ is the
    downside deviation. This metric rewards agents that avoid left-tail
    losses rather than simply minimising total variance.

    \item \textbf{Mean Total Fees}: Cumulative transaction costs in wealth
    units, averaged across simulation paths.
\end{itemize}

\subsection{Statistical Tests}

\begin{itemize}
    \item \textbf{Fisher's Exact Test}: Applied to the $2 \times 2$
    contingency table of ruin / survival counts for two agents. Tests the
    null hypothesis $H_0: p_{\text{ruin}}^{(A)} = p_{\text{ruin}}^{(B)}$
    against the two-sided alternative.

    \item \textbf{Mann-Whitney U Test}: Applied to the terminal wealth
    distributions of two agents. Non-parametric; appropriate given the
    log-normal / power-law shape of compounding wealth. Tests
    $H_0: \Prob(W_T^{(A)} > W_T^{(B)}) = 1/2$.
\end{itemize}

We report $p$-values and reject $H_0$ at the $\alpha = 0.05$ significance
level throughout.
